\clearpage
\fancyhead[L]{\small GAUSS TO BODE.}
\hypertarget{native-289-lower}{}
\begin{flushright}
Brunswick, 23 August 1807.
\end{flushright}

\par\indent Please forgive me for sending you no further news of my later work on this subject since I sent you my first elements of Vesta.
\SourceLeaf{290}%
\par\indent A trip to Bremen this summer and an illness I suffered afterward set me back somewhat in my correspondence as well as in my work generally. I am now all the more eager to make up for the delay, so that you can still make whatever use you wish of the information I send you in your new Yearbook.

\par\indent From the \href{https://zs.thulb.uni-jena.de/rsc/viewer/jportal_derivate_00237836/Monatliche_Correspondenz_130168688_15_1807_0605.tif}{M.\,C.} you will already know my second set of elements for \VestaBodyTwoNinety{} and their agreement with a great many observations.

\par\indent To calculate the \emph{third} set of elements, I used the Paris and Milan meridian observations and the latest observations from Bremen and Lilienthal --- which extend to 11 July --- in addition to the earlier observations. Here is this third set of elements.

\begin{center}\begingroup\fontsize{9}{11}\selectfont
\begin{tabular}{ll}
\shortstack[l]{Epoch of longitude, 31 March 1807\\$0^{\text{u}}$ in Bremen} .\ .\ . & $192^\circ\,23'\,30''1$ \\
Aphelion .\ .\ .\ .\ .\ .\ .\ .\ .\ .\ .\ .\ .\ .\ .\ . & $69\;50\;31.9$ \\
Node .\ .\ .\ .\ .\ .\ .\ .\ .\ .\ .\ .\ .\ .\ .\ .\ .\ . & $103\;18\;28.0$
\end{tabular}\endgroup\end{center}

\noindent Both for the above epoch and assumed to be fixed relative to the stars.

\begin{center}\begingroup\fontsize{9}{11}\selectfont
\begin{tabular}{ll}
Inclination of the orbit .\ .\ .\ .\ .\ .\ .\ .\ .\ .\ .\ .\ . & $7\;8\;10.7$ \\
Eccentricity .\ .\ .\ .\ .\ .\ .\ .\ .\ .\ .\ .\ .\ .\ . & $0.0855050$ \\
Logarithm of the semiaxis .\ .\ .\ .\ .\ .\ .\ .\ .\ . & $0.3720160$ \\
Mean daily tropical motion .\ .\ .\ .\ .\ . & $981''7087$
\end{tabular}\endgroup\end{center}

\par\indent I also enclose the ephemeris for next year, calculated from these elements.

[Right ascension and declination are given to minute precision for midnight at Seeberg at four-day intervals from 2 May 1808 to 18 March 1809.]

